\chapter{Inequalities}

\begin{center}\itshape
The number $e$ appears in an inequality about averages of an arbitrary positive sequence.\\
Why there, and why is it the best possible constant?
\end{center}

\section{Convexity and Jensen}

\begin{reminder}[convex functions]
$f:I\to\R$ is \emph{convex} if $f(\lambda x+(1-\lambda)y)\le\lambda f(x)+(1-\lambda)f(y)$ for $x,y\in I$, $\lambda\in[0,1]$:
chords lie above the graph. If $f$ is twice differentiable, $f$ is convex if and only if $f''\ge0$.
A convex function lies above its tangents. \emph{Concave} means $-f$ convex.
\end{reminder}

\begin{theorem}[Jensen]\label{thm:jensen}
Let $f$ be convex on $I$, $x_1,\dots,x_n\in I$ and $\lambda_1,\dots,\lambda_n\ge0$ with $\sum\lambda_i=1$. Then
\[
f\left(\sum_{i=1}^n\lambda_ix_i\right)\le\sum_{i=1}^n\lambda_if(x_i).
\]
\end{theorem}
\begin{proof}
Induction on $n$, the case $n=2$ being the definition. If $\lambda_n=1$ there is nothing to prove. Otherwise put
$\mu=1-\lambda_n$ and $y=\sum_{i<n}\frac{\lambda_i}\mu x_i\in I$. Then $f\left(\mu y+\lambda_nx_n\right)\le\mu f(y)+\lambda_nf(x_n)$, and
$f(y)\le\sum_{i<n}\frac{\lambda_i}\mu f(x_i)$ by the induction hypothesis.
\end{proof}

\begin{theorem}[AM--GM]\label{thm:amgm}
For $x_1,\dots,x_n>0$ and weights $\lambda_i\ge0$ with $\sum\lambda_i=1$,
\[
\prod_{i=1}^nx_i^{\lambda_i}\le\sum_{i=1}^n\lambda_ix_i,\qquad\text{in particular}\qquad\sqrt[n]{x_1\cdots x_n}\le\frac{x_1+\dots+x_n}n,
\]
with equality in the unweighted case if and only if all the $x_i$ are equal.
\end{theorem}
\begin{proof}
Apply Jensen to the concave function $\ln$ and take exponentials. Since $\ln$ is strictly concave, equality forces all the $x_i$ to be equal.
\end{proof}

\section{Cauchy--Schwarz and H\"older}

\begin{theorem}[Cauchy--Schwarz]\label{thm:cs}
For real numbers $a_i,b_i$,
$\displaystyle\left(\sum_{i=1}^na_ib_i\right)^2\le\left(\sum_{i=1}^na_i^2\right)\left(\sum_{i=1}^nb_i^2\right)$,
with equality if and only if the vectors $a$ and $b$ are proportional.
\end{theorem}
\begin{proof}
The polynomial $t\mapsto\sum_i(a_i+tb_i)^2=\sum a_i^2+2t\sum a_ib_i+t^2\sum b_i^2$ is non-negative, so if $b\ne0$ its discriminant is $\le0$, which is the inequality.
Equality means a double root $t$, that is $a+tb=0$.
\end{proof}

The same proof works in any real inner product space: $\langle a,b\rangle^2\le\langle a,a\rangle\langle b,b\rangle$. For instance, for continuous $f,g$ on $[a,b]$,
$\left(\int fg\right)^2\le\int f^2\int g^2$.

A frequently used form (``Engel'' or ``Titu''): for $x_i\in\R$ and $y_i>0$,
\[
\sum_{i=1}^n\frac{x_i^2}{y_i}\ge\frac{\left(\sum x_i\right)^2}{\sum y_i},
\]
obtained from Cauchy--Schwarz with $a_i=x_i/\sqrt{y_i}$ and $b_i=\sqrt{y_i}$.

\begin{reminder}[rearrangement]
If $a_1\le\dots\le a_n$ and $b_1\le\dots\le b_n$, then for every permutation $\sigma$,
$\sum a_ib_{n+1-i}\le\sum a_ib_{\sigma(i)}\le\sum a_ib_i$: a sum of products is largest when both lists are sorted the same way.
(Proof: if $\sigma$ is not increasing, swapping two values $i<j$ with $\sigma(i)>\sigma(j)$ changes the sum by $(a_j-a_i)(b_{\sigma(i)}-b_{\sigma(j)})\ge0$.)
\end{reminder}

\section{Exercises}

\begin{exercise}[$\star$]\label{ex:8abc}
Show that $(a+b)(b+c)(c+a)\ge8abc$ for $a,b,c>0$.
\end{exercise}

\begin{exercise}[$\star$]\label{ex:sqrtk}
Show that $\displaystyle\sum_{k=1}^n\sqrt k\le n\sqrt{\frac{n+1}2}$.
\end{exercise}

\begin{exercise}[$\star$]\label{ex:intcs}
Let $f$ be continuous and positive on $[0,1]$. Show that $\displaystyle\int_0^1f\cdot\int_0^1\frac1f\ge1$. When is there equality?
\end{exercise}

\begin{exercise}[$\star\star$]\label{ex:euler}
Show that the sequence $\left(1+\frac1n\right)^n$ is strictly increasing. \emph{Hint: AM--GM on $n+1$ well-chosen numbers.}
\end{exercise}

\begin{exercise}[$\star\star$]\label{ex:nesbitt}
\emph{Nesbitt's inequality.} Show that $\dfrac a{b+c}+\dfrac b{c+a}+\dfrac c{a+b}\ge\dfrac32$ for $a,b,c>0$.
\end{exercise}

\begin{exercise}[$\star\star$]\label{ex:triangle}
Show that in every triangle with angles $A,B,C$, $\sin A+\sin B+\sin C\le\dfrac{3\sqrt3}2$.
\end{exercise}

\begin{exercise}[$\star\star$]\label{ex:aabbcc}
Show that $a^ab^bc^c\ge(abc)^{(a+b+c)/3}$ for $a,b,c>0$. \emph{Hint: $(x-y)(\ln x-\ln y)\ge0$.}
\end{exercise}

\begin{exercise}[$\star\star$]\label{ex:holder}
\begin{enumerate}[label=(\alph*)]
\item \emph{Young.} Let $p,q>1$ with $\frac1p+\frac1q=1$. Show that $ab\le\frac{a^p}p+\frac{b^q}q$ for $a,b\ge0$.
\item \emph{H\"older.} Deduce that $\sum a_ib_i\le\left(\sum a_i^p\right)^{1/p}\left(\sum b_i^q\right)^{1/q}$ for $a_i,b_i\ge0$.
\end{enumerate}
\end{exercise}

\begin{exercise}[$\star\star\star$]\label{ex:carleman}
\emph{Carleman's inequality.} Let $(a_n)_{n\ge1}$ be non-negative with $\sum a_n<\infty$. Show that
\[
\sum_{n=1}^\infty\left(a_1a_2\cdots a_n\right)^{1/n}\le e\sum_{n=1}^\infty a_n.
\]
\emph{Hint: with $c_k=\frac{(k+1)^k}{k^{k-1}}$, compute $c_1c_2\cdots c_n$ and apply AM--GM to $c_1a_1,\dots,c_na_n$.}
\end{exercise}

\begin{beauty}
The answer to the opening question is in Exercise~\ref{ex:carleman}: $e$ enters through $\left(1+\frac1k\right)^k<e$, the sequence of
Exercise~\ref{ex:euler}. The notebook shows numerically that no smaller constant works. In Tome~3, Jensen becomes
$f(\mathbb E[X])\le\mathbb E[f(X)]$, and Cauchy--Schwarz bounds covariances; in Tome~4, H\"older is the key to $L^p$ spaces.
\end{beauty}

\section{Solutions}

\begin{solution}{ex:8abc}
By AM--GM, $a+b\ge2\sqrt{ab}$, $b+c\ge2\sqrt{bc}$, $c+a\ge2\sqrt{ca}$. Multiply: the product is at least $8\sqrt{a^2b^2c^2}=8abc$.
\end{solution}

\begin{solution}{ex:sqrtk}
Cauchy--Schwarz with $a_k=1$, $b_k=\sqrt k$: $\left(\sum\sqrt k\right)^2\le n\sum k=\frac{n^2(n+1)}2$. Take square roots.
\end{solution}

\begin{solution}{ex:intcs}
Cauchy--Schwarz for integrals with $\sqrt f$ and $1/\sqrt f$: $1=\left(\int\sqrt f\cdot\frac1{\sqrt f}\right)^2\le\int f\int\frac1f$.
Equality holds if and only if $\sqrt f$ and $1/\sqrt f$ are proportional, that is, $f$ is constant.
\end{solution}

\begin{solution}{ex:euler}
Apply AM--GM to the $n+1$ numbers $1,1+\frac1n,\dots,1+\frac1n$ ($n$ copies). Their mean is $\frac{1+(n+1)}{n+1}=1+\frac1{n+1}$, so
\[
\left(1+\frac1n\right)^{n/(n+1)}<1+\frac1{n+1},
\]
strictly since the numbers are not all equal. Raising to the power $n+1$ gives $\left(1+\frac1n\right)^n<\left(1+\frac1{n+1}\right)^{n+1}$.
\end{solution}

\begin{solution}{ex:nesbitt}
Let $S=a+b+c$. Then $\frac a{b+c}=\frac S{S-a}-1$, so the left-hand side is $S\sum\frac1{S-a}-3$. By the Engel form of Cauchy--Schwarz
with $x_i=1$,
\[
\sum\frac1{S-a}\ge\frac9{(S-a)+(S-b)+(S-c)}=\frac9{2S}.
\]
Hence the sum is at least $\frac92-3=\frac32$, with equality if and only if $a=b=c$.
\end{solution}

\begin{solution}{ex:triangle}
$\sin$ is concave on $[0,\pi]$ since $\sin''=-\sin\le0$ there. By Jensen with weights $\frac13$,
$\frac{\sin A+\sin B+\sin C}3\le\sin\frac{A+B+C}3=\sin\frac\pi3=\frac{\sqrt3}2$. Equality for the equilateral triangle.
\end{solution}

\begin{solution}{ex:aabbcc}
$\ln$ is increasing, so $(x-y)(\ln x-\ln y)\ge0$. Sum over the three pairs $\{a,b\},\{b,c\},\{c,a\}$:
\[
2(a\ln a+b\ln b+c\ln c)\ge\sum_{x\ne y}x\ln y,
\]
where $x,y$ range over $a,b,c$. Adding $a\ln a+b\ln b+c\ln c$ to both sides gives $3\sum a\ln a\ge\sum_{x,y}x\ln y=(a+b+c)(\ln a+\ln b+\ln c)$;
divide by $3$ and exponentiate.
\end{solution}

\begin{solution}{ex:holder}
(a) If $a$ or $b$ is $0$ it is clear. Otherwise, by weighted AM--GM with weights $\frac1p,\frac1q$ applied to $a^p$ and $b^q$:
$ab=(a^p)^{1/p}(b^q)^{1/q}\le\frac{a^p}p+\frac{b^q}q$.

(b) Let $A=\left(\sum a_i^p\right)^{1/p}$ and $B=\left(\sum b_i^q\right)^{1/q}$; if $A$ or $B$ is $0$ it is clear. Apply (a) to $a_i/A$ and $b_i/B$ and sum:
\[
\frac{\sum a_ib_i}{AB}\le\frac1p\frac{\sum a_i^p}{A^p}+\frac1q\frac{\sum b_i^q}{B^q}=\frac1p+\frac1q=1.
\]
\end{solution}

\begin{solution}{ex:carleman}
The product telescopes: $\prod_{k=1}^nc_k=\dfrac{2^1\,3^2\cdots(n+1)^n}{1^0\,2^1\cdots n^{n-1}}=(n+1)^n$. So, by AM--GM,
\[
\left(a_1\cdots a_n\right)^{1/n}=\frac{\left(c_1a_1\cdots c_na_n\right)^{1/n}}{n+1}\le\frac1{n(n+1)}\sum_{k=1}^nc_ka_k.
\]
Sum over $n\ge1$ and swap the sums (all terms are non-negative), using $\sum_{n\ge k}\frac1{n(n+1)}=\frac1k$ (telescoping):
\[
\sum_{n=1}^\infty\left(a_1\cdots a_n\right)^{1/n}\le\sum_{k=1}^\infty\frac{c_k}ka_k=\sum_{k=1}^\infty\left(1+\frac1k\right)^ka_k\le e\sum_{k=1}^\infty a_k,
\]
since $\left(1+\frac1k\right)^k$ increases to $e$ (Exercise~\ref{ex:euler}).
\end{solution}
